% TOP_PROBLEM: {"id": 30002415, "title": "Sharp Multidimensional Star-Discrepancy Asymptotics", "category_id": 8, "category": {"id": 8, "name": "analysis", "display_name": "Analysis", "description": "Limits, continuity, calculus, and function theory.", "slug": "analysis", "order_index": 8, "created_at": "2026-07-31T15:26:25.671Z"}, "status": "open"}
% FIRST_POSTED: 2026-09-27
% !TeX program = lualatex
% ATTEMPT_STATUS: unresolved
% Model: GPT-6 Astra Ultra; continuation dated 27 September 2026.
\documentclass[11pt]{article}
\usepackage[margin=25mm]{geometry}
\usepackage{fontspec}
\setmainfont{Latin Modern Roman}
\usepackage{amsmath,amssymb,mathtools}
\usepackage{unicode-math}
\setmathfont{Latin Modern Math}
\usepackage{xurl,hyperref}
\hypersetup{colorlinks=true,urlcolor=blue}
\setcounter{secnumdepth}{0}
\setlength{\parindent}{0pt}
\setlength{\parskip}{5pt}
\setlength{\emergencystretch}{3em}
\begin{document}
\section*{\texorpdfstring{Sharp Multidimensional Star-Discrepancy Asymptotics}{Sharp Multidimensional Star-Discrepancy Asymptotics}}\label{30002415}
\subsection{Definitions and mathematical statement}
For an $N$-point set $P\subset[0,1)^d$, its unnormalized star discrepancy is
\[
 D^*(P)=\sup_{x\in[0,1]^d}
 \left|\#\left(P\cap\prod_{j=1}^d[0,x_j)\right)-N\prod_{j=1}^dx_j\right|.
\]
Let $\Delta_d(N)=\inf_{|P|=N}D^*(P)$. For every fixed $d\ge3$, determine a function $f_d(N)$ for which $c_df_d(N)\le\Delta_d(N)\le C_df_d(N)$ for all sufficiently large $N$, with positive constants depending only on $d$. The boxes are anchored at the origin and half-open. This convention omits the factor $1/N$ used in normalized discrepancy, so any proposed asymptotic must match this normalization. The dimension is fixed while the number of points grows.
\par\smallskip\textit{Formulation and concept references:} \href{https://www-users.cse.umn.edu/~dbilyk/bilyk-research.pdf}{[S1]}.

\subsection{Short English statement}
How small can the worst counting error in origin-anchored boxes be for an optimally placed $N$-point set in each fixed dimension at least three?
\subsection{Research attempt}
\paragraph{Normalization and the unresolved comparison.}
The discrepancy here is the counting error itself. We write
\[
 D_P(x)=\sum_{p\in P}\mathbf1_{\{p_j<x_j\ \forall j\}}
                         -N\prod_{j=1}^dx_j.
\]
The primary source [S1, Section 1.1] uses this normalization and
explains the gap between known supremum-norm estimates in dimension
at least three. Results on inverse discrepancy when dimension varies
ask a different asymptotic question. A normalized discrepancy bound
must be multiplied by $N$ before comparison with this entry.

We give an orthogonal-function lower bound with explicit constants,
an elementary existence upper bound, exact finite evaluation rules,
and a product-grid stress test. The bounds proved here do not match;
their remaining gap is stated quantitatively rather than hidden in
the phrase that the points are uniformly distributed.

\paragraph{An elementary one-dimensional obstruction.}
In dimension one, order the points $p_1<\cdots<p_N$. At $x=p_j$
the half-open counting function has value $j-1$, while its limit
from the right is $j$. Thus the discrepancy has two one-sided values
whose difference is one. At least one has absolute value at least
$1/2$, giving $D^*(P)\geq1/2$. The midpoint set
$p_j=(j-1/2)/N$ attains this value, including both one-sided limits.
Hence $\Delta_1(N)=1/2$ in the unnormalized convention.

Projection onto any coordinate gives the same lower bound in higher
dimensions, since the remaining box coordinates can be set equal
to one. Much stronger growth is forced by multiscale geometry, as
the following complete argument shows.

\paragraph{Dyadic Haar functions and empty rectangles.}
For a dyadic interval $I=[a,a+\ell)$, let $h_I$ be minus one on
its left half, plus one on its right half, and zero outside $I$.
Then
\[
 \int h_I=0,\qquad \|h_I\|_2^2=\ell,\qquad
 \int_0^1 xh_I(x)\,dx=\ell^2/4.
 \tag{368.1}
\]
For a dyadic rectangle $R=I_1\times\cdots\times I_d$, set
$h_R(x)=\prod_jh_{I_j}(x_j)$. Its squared norm is $|R|$.

Suppose the interior of $R$ contains no point of $P$. For one
point $p$, its contribution to $\langle D_P,h_R\rangle$ factors
into $d$ integrals of the form
$\int_{I_j}\mathbf1_{p_j<x_j}h_{I_j}(x_j)\,dx_j$.
Unless every $p_j$ is strictly inside $I_j$, one factor vanishes:
the indicator is either zero or constant one almost everywhere on
that interval. Points on its endpoints cause no exception. Thus all
point contributions vanish for an empty-interior rectangle. The
volume term and (368.1) give exactly
\[
 \langle D_P,h_R\rangle=-\frac{N|R|^2}{4^d}.
 \tag{368.2}
\]
The sign is irrelevant for the norm estimate but fixes the convention
for checking the calculation.

Choose an integer $m$ with $2N\leq2^m<4N$. For each vector of
nonnegative integers $r=(r_1,\ldots,r_d)$ with $\sum_jr_j=m$,
the rectangles with side lengths $2^{-r_j}$ partition the unit cube
into $2^m$ pieces. At most $N$ have an interior point. Hence at
least $2^m-N\geq2^{m-1}$ are empty in the required sense.

All Haar functions selected this way, across all shapes $r$, are
orthogonal. In one dimension, two distinct dyadic Haar functions
have disjoint supports or the smaller support lies in a half on
which the larger is constant, and the smaller has mean zero.
Taking tensor products proves orthogonality in $d$ dimensions as
soon as the rectangles differ in some coordinate interval.

\paragraph{A complete logarithmic lower bound.}
Bessel's inequality applied to the normalized empty-rectangle Haar
functions gives
\[
 \|D_P\|_2^2
 \geq\sum_{r_1+\cdots+r_d=m}\ \sum_{R\text{ empty of shape }r}
       \frac{|\langle D_P,h_R\rangle|^2}{|R|}.
\]
Every summand is $N^2 2^{-3m}/16^d$, and there are at least
$2^{m-1}$ per shape. The number of shapes is
$\binom{m+d-1}{d-1}$, by the elementary stars-and-bars count.
Using $2^m<4N$ yields
\[
 \|D_P\|_2^2\geq
       \frac1{32\,16^d}\binom{m+d-1}{d-1}.
 \tag{368.3}
\]
Since the cube has volume one, $\|D_P\|_\infty\geq\|D_P\|_2$.
Consequently, for fixed $d$ and $N\geq2$,
\[
 D^*(P)\geq c_d(\log N)^{(d-1)/2}.
 \tag{368.4}
\]
The use of an essential supremum in the norm does not weaken this
conclusion for the stated supremum over boxes. The latter is at
least the former. No assumption that the points avoid dyadic
boundaries was made; the empty-interior argument handles them.

This is the classical orthogonal-function mechanism, reproduced
here rather than merely quoted. It does not give the conjecturally
sharp supremum exponent. The square root in Bessel's inequality
combines the contributions of the different rectangle shapes in
an $\ell^2$ fashion. Replacing that square root by an ordinary sum
would require a new estimate and is not an algebraic simplification.

\paragraph{A direct existence upper bound for every number of points.}
Let $X_1,\ldots,X_N$ be independent uniform points in the cube.
Take the grid $G=\{0,1/N,\ldots,1\}^d$. For each $x\in G$,
the box count is a sum of independent Bernoulli variables with
mean $N\prod_jx_j$. The elementary exponential-moment estimate
for a Bernoulli variable gives Hoeffding's bound
\[
 \mathbb P(|D_P(x)|>t)\leq2e^{-2t^2/N}.
 \tag{368.5}
\]
One derivation centers each summand, bounds its moment generating
function by $e^{\lambda^2/8}$, multiplies by independence and
optimizes Markov's inequality at $\lambda=4t/N$; the lower tail
uses $-\lambda$.

The union bound over $(N+1)^d$ corners shows that with positive
probability all grid discrepancies have absolute value at most
\[
 t_N=\sqrt{\frac N2\bigl(d\log(N+1)+\log4\bigr)}.
 \tag{368.6}
\]
For an arbitrary $x$, let $x^-$ and $x^+$ be its coordinatewise
lower and upper grid neighbors. Box counts are monotone, and
\[
 0\leq\prod_jx_j^+-\prod_jx_j^-\leq\sum_j(x_j^+-x_j^-)\leq d/N.
\]
Sandwiching the count and subtracting the true volume therefore
gives $|D_P(x)|\leq t_N+d$. Random points are distinct and lie
in $[0,1)^d$ almost surely. We have proved, for every $N$,
\[
 \Delta_d(N)\leq
 \sqrt{\frac N2\bigl(d\log(N+1)+\log4\bigr)}+d.
 \tag{368.7}
\]
This bound is deliberately elementary and is weaker than the
specialized low-discrepancy constructions discussed in [S1]. It
is not presented as the best known upper bound. Its role here is
to supply a complete existence argument with the same normalization
as the lower bound.

\paragraph{Why a regular product grid is not automatically near-optimal.}
For $N=q^d$, take the product of one-dimensional midpoint grids.
Let $a_j(x_j)$ be the number of coordinate midpoints below $x_j$.
Then $|a_j-qx_j|\leq1/2$, interpreted by the appropriate one-sided
limit at a midpoint, and both $a_j$ and $qx_j$ lie in $[0,q]$.
A product telescoping identity gives
\[
 \left|\prod_ja_j-q^d\prod_jx_j\right|
      \leq\frac d2 q^{d-1}.
 \tag{368.8}
\]
There is also a matching obstruction in order: choose $x_1$ just
above the first midpoint and set the other coordinates to one.
The discrepancy tends to $q^{d-1}/2$. Thus these attractive
geometric grids have discrepancy of order $N^{1-1/d}$ for fixed
$d$, far above logarithmic scales. Small nearest-neighbor spacing
and visual regularity do not certify the anchored-box objective.

\paragraph{An exact finite formula respecting half-open boxes.}
For computation, let
$\Gamma_j=\{0,1\}\cup\{p_j:p\in P\}$ and
$\Gamma=\prod_j\Gamma_j$. At each corner define
\[
 A(x)=\#\{p:p_j<x_j\ \forall j\},\qquad
 \overline A(x)=\#\{p:p_j\leq x_j\ \forall j\}.
\]
Then the exact value is
\[
 D^*(P)=\max_{x\in\Gamma}
 \max\left\{N\prod_jx_j-A(x),\
                   \overline A(x)-N\prod_jx_j\right\}.
 \tag{368.9}
\]
To prove it, inside any open cell cut out by the coordinate grids
the point count is constant, while the volume product is monotone
in every coordinate. The most positive and most negative errors
are therefore attained as one-sided limits at its corners. At an
upper limiting corner, the unchanged half-open count is $A$; at a
lower limiting corner, approaching from above includes equality
points and gives $\overline A$. Corners at zero are obtained by
positive limiting coordinates, and no data point has coordinate
one, so the boundary at one causes no exception. Conversely each
term in (368.9) is an actual value or such a one-sided limiting
value of the discrepancy, and hence is bounded by its supremum.

Ties among point coordinates are handled by counting all equality
points together. Evaluating only the half-open counts at data
corners misses the second family of limits and can underestimate
the discrepancy. Formula (368.9) provides a finite certificate for
one supplied point set, not a proof that it minimizes discrepancy
among all real configurations.

\paragraph{What would strengthen the lower-bound route.}
The Haar coefficients in (368.2) live on many different rectangle
shapes of the same volume. They are orthogonal, but their pointwise
signs overlap in complicated ways. One may attempt a nonlinear test
function built from products of shape sums, paired against $D_P$.
To obtain a sharper $L^\infty$ lower bound, its $L^1$ norm must be
controlled while its correlation with the discrepancy accumulates
more efficiently than (368.3). Cross terms from rectangles sharing
coordinates are then essential. Ignoring those terms is not justified
by orthogonality, which controls pairwise integrals rather than
all products arising in such a construction.

The source discusses stronger specialized lower bounds and the
small-ball approach; this attempt does not reproduce those results
or turn them into a sharp all-dimensional order. The elementary
proof isolates why simple square-function information alone stops
at the exponent in (368.4).

\paragraph{Diagnostics and outcome.}
The script evaluates (368.9) exactly for rational test point sets,
checks the midpoint-grid one-dimensional value and the product-grid
obstruction, and verifies the empty-rectangle Haar integral by
piecewise polynomial integration. These tests audit endpoint
conventions and constants. Sampling random boxes alone would not
be an exact supremum computation, so it is not used as a certificate.

The proved lower and upper bounds, (368.4) and (368.7), do not match.
The first unresolved step in the proposed harmonic route is a
uniform nonlinear combination of shape contributions strong enough
to reach a matching supremum estimate without uncontrolled cross
terms. Determining $f_d(N)$ in the printed two-sided sense remains
unresolved by this attempt.
\subsection{Sources}
\begin{itemize}
\item[S1]Dmitriy Bilyk, Research Statement, Section 1.1.
\url{https://www-users.cse.umn.edu/~dbilyk/bilyk-research.pdf}.
\textit{Formulation source.}
\end{itemize}
\end{document}
